\section{실현착순 외부검증}
\label{sec:external-validation}

마지막으로 내부 가격정합성과 분리하여, clean Panel A의 3,321개 경주에서
삼쌍승 주변가격과 단승--Harville 가격측도가 실제 착순에 부여한 로그점수를
비교한다. 각 목표 승식의 \texttt{is\_hit} 적중조합이 정확히 하나인지, 파싱한
상위 3착이 그 조합과 일치하는지 전수검사한다. 이 검사는 동착으로 복수
적중조합이 생기거나 착순 파싱이 어긋난 경주를 점로그점수로 조용히 처리하지
않도록 한다. 현재 clean 표본에서는 두 조건을 위반한 경주가 0개다.

\begin{table}[htbp]
\centering
\caption{Clean Panel A의 실현착순 로그점수 개선. 개선폭은 단승--Harville
로그점수에서 삼쌍승 주변화 로그점수를 뺀 값으로, 양수일수록 삼쌍승 주변화가
낫다. 미보정 평균은 경주일 군집 bootstrap, 미보정 중앙값은 경주 bootstrap
95\% 구간을 함께 보고한다. 보정 열은 두 가격측도에 같은 연도제외 단조보정을
각각 적용한 결과다. EPS 구속은 `미보정 main/Harville; 보정 main/Harville'
순서의 $10^{-12}$ 하한 구속 관측 수다. 단승 행도 항등식이 아니라 삼쌍승 주변 단승가격과 실제 단승가격의
실질 비교다.}
\label{tab:main-external-log-scores}
\resizebox{\textwidth}{!}{\begin{tabular}{lrrrrr}
\toprule
승식 & 경주 수 & 삼쌍승 주변화 & 단승--Harville & 개선폭 & 95\% CI \\
\midrule
단승 & 3,321 & 1.8350 & 1.8345 & -0.0005 & [-0.0064, 0.0064] \\
쌍승 & 3,321 & 3.6583 & 3.6900 & 0.0317 & [0.0197, 0.0444] \\
복승 & 3,321 & 2.9903 & 3.0188 & 0.0285 & [0.0169, 0.0388] \\
삼복승 & 3,321 & 3.7016 & 3.7776 & 0.0759 & [0.0602, 0.0920] \\
\bottomrule
\end{tabular}
}
\end{table}

로그점수는 기대손실이므로 경주균등 산술평균을 주집계치로 사용하고, 같은 날의
공통충격을 허용하도록 경주일을 재표집한다. 이는 주패널의 사전 집계치인 중앙값을
대체하지 않는 별도 외부검증이다. 사전 규칙과 직접 비교할 수 있도록 경주별
개선폭 중앙값과 경주 bootstrap 구간도 병기한다. 또한 각 승식·모형에서 다른
연도를 훈련자료로 삼아 예측가격과 실현지시자의 단조관계를 추정하고, 검증연도
각 경주 안에서 다시 합이 1이 되도록 정규화한 대칭적 보정 결과를 제시한다.

미보정 가격측도의 평균 개선은 쌍승 0.0317, 복승 0.0285, 삼복승 0.0759로
양수이고 단승은 $-0.0005$로 구분되지 않는다. 경주별 중앙값과 경주 bootstrap도
세 복합승식에서 양수다. 대칭적 연도제외 보정 뒤 평균 개선은 각각 0.0275,
0.0235, 0.0423으로 줄지만 경주일 군집 구간이 여전히 0보다 크며, 단승은
0.0051로 구분되지 않는다. 미보정에서는 확률하한 구속이 0건이고, 보정 뒤에는
세 복합승식에서 두 모형 각각 1건이 구속된다. 이 드문 끝점 관측을 포함해도
복합승식의 방향은 유지된다.

그러나 미보정 값은 객관확률이 아니라 공제와 인기편향이 섞인 두 가격측도의
상대 왜곡을 비교한다. 따라서 핵심 해석은 보정 전후의 부호와 민감도에 한정한다. 표시상한 조합의 잠재
로그점수는 점식별되지 않으므로 이 결과를 상한 경주, 인과적 정보량 우위 또는
객관적 시장효율성으로 일반화하지 않는다.
